\section{Countermodels and non-substitution laws}
\label{sec:countermodels}

The following finite examples separate the exact conclusions of the
duality theory from the physical estimates that remain open.

\subsection{A shared rule creates a real joint gap}

Let one scalar \(t\in\mathbb R\) control two edge completions, with costs
\[
 p_{12}(t)=|t|,
 \qquad
 p_{23}(t)=|1-t|.
\]
Each edge separately has optimal cost zero.  Under the shared rule,
\[
 \inf_t\max\{p_{12}(t),p_{23}(t)\}
 =
 \frac12.
\]
Thus a joint gap can be real, but it comes from the common parameter,
not from Hermitian reversal.

\subsection{Adaptive choice reverses max and min}

In the same example,
\[
 \max_{e\in\{12,23\}}\inf_t p_e(t)=0,
 \qquad
 \inf_t\max_{e\in\{12,23\}}p_e(t)=\frac12.
\]
Choosing the completion after the worst edge or operator extremizer has
been revealed changes the problem.  A coefficient-uniform theorem must
use the second order of quantifiers.

\subsection{Fixed-pair decay does not control a matrix}

Let \(P_N\) be the symmetric zero-diagonal matrix with
\[
 (P_N)_{mn}=\frac1{N-1}\qquad(m\ne n).
\]
Every fixed entry tends to zero, but
\[
 \rho(P_N)=1.
\]
Hence even optimal entrywise completion costs of order \(o(1)\) are
insufficient without a row, weighted-Schur, or Perron aggregation over
the complete growing active set.

\subsection{The absolute Perron floor can miss signed cancellation}

Let \(N=2^r\), and let \(S_N\) be the symmetric Sylvester Hadamard matrix.
Then
\[
 S_N=S_N^{\mathsf T},
 \qquad
 S_NS_N^{\mathsf T}=NI,
 \qquad
 |(S_N)_{mn}|=1.
\]
Define the real symmetric zero-diagonal matrix
\[
 H_N=\frac{S_N-\diag(S_N)}{N}.
\]
Since \(\|S_N\|=\sqrt N\),
\[
 \|H_N\|
 \le
 \frac1{\sqrt N}+\frac1N
 \longrightarrow0.
\]
On the other hand,
\[
 |H_N|=\frac{J-I}{N},
 \qquad
 \rho(|H_N|)=\frac{N-1}{N}\longrightarrow1.
\]
By \cref{thm:Perron-floor}, every entrywise absolute completion majorant
of this signed matrix has spectral radius asymptotically at least one.
Thus failure of the whole completion--Perron architecture need not be
failure of signed operator cancellation.

\subsection{One bad completion is not an obstruction}

Take zero primitive datum, a prime \(p>D\), and a rectangle containing
the nonprimitive site \((p,p)\).  The fill
\[
 F_A=Ae_{p,p}
\]
is legal for every \(A\in\C\).  The tail coordinate
\((d,x,y)=(p,1,1)\) has unit weight, so its exact tail cost is at least
\(|A|\).  The optimum remains zero because the zero completion is legal.
A large cost for a chosen interpolation proves nothing about
\(\gamma_D(0)\); only an optimal dual witness can lower-bound the
infimum.

\subsection{A dual obstruction is not a signed lower bound}

In the abstract normal form \(R=I\), \(A=I\), take
\[
 k=(1,-1),
 \qquad
 \sigma=(1,1).
\]
Then the signed functional vanishes:
\[
 \langle\sigma,k\rangle=0,
\]
whereas the completion gauge is
\[
 \gamma(k)=\|k\|_1=2.
\]
The optimal dual witness obstructs a cheap \(\ell^1\) certificate, not
the signed functional.  The Hadamard construction above is the matrix
version of the same separation.

\subsection{Costs before content aggregation can discard cancellation}

Suppose two literal content slices satisfy
\[
 L_{mn,q_1}=F,
 \qquad
 L_{mn,q_2}=-F.
\]
The physically relevant post-\(q\)-aggregation datum is zero.  Completing
the two slices separately and adding their absolute costs can be
positive, whereas completing their exact sum costs zero.  Therefore a
paper must declare whether its completion rule acts after exact
\(q\)-aggregation or retains \(q\) as a coupled label.  One may not move
between the two orders after seeing which is cheaper.

\subsection{Sidewise renormalization changes the object}

Suppose two orthogonal sides each contribute row mass one, so the raw
global masses are \(W_m=W_n=2\).  If one side contains one common unit
key, its correct normalized two-corner contribution is
\[
 \frac1{\sqrt{W_mW_n}}=\frac12.
\]
Renormalizing that side alone gives \(1\), doubling the physical entry.
Analytic fill values never enter \(W_m\) or \(W_n\); a completion cannot
authorize a new normalization.

\subsection{Hermitian-incompatible scalar completions}

For zero primitive data, both
\[
 F_{12}=Ae_{2,2},
 \qquad
 F_{21}=0
\]
are legal scalar completions.  Unless \(A=0\), they do not satisfy
\[
 F_{21}(v,u)=\overline{F_{12}(u,v)}.
\]
They cannot be inserted together into one Hermitian analytic carrier.
One must use upper-triangle variables or apply
\cref{prop:Hermitian-symmetrization}.

\subsection{A fixed weight is not automatically Perron-optimal}

For a nonnegative matrix \(P\) and an arbitrary \(c>0\),
\[
 \rho(P)\le\max_m\frac{(Pc)_m}{c_m},
\]
and the inequality may be strict.  Equality requires an appropriate
Perron weight or the infimum over all positive \(c\).  A dual lower
witness for one fixed-\(c\) Schur problem does not by itself lower-bound
the infimum over \(c\).
